\section{Finite-Observation Transfer Proof}\label{app:r14sensingproof}
\subsection{Proof of Proposition~\ref{prop:r14sensor}}
All processes are coupled using the same economic Brownian motion and
initial profile; the sensor noise uses the independent randomizer.
The real-map Lipschitz bound, uniform implementation error, and
nonexpansiveness of the common guard give
\[
 \|\widehat\pi_k(\overline Y_{t_k}+\eta_k)
                  -\widehat\pi_k(Z_k)\|_{2,d}
 \leq L\{\|\overline Y_{t_k}-Z_k\|_{2,d}+\nu_{\rm obs}\}+2\delta_\pi.
\]
The common action tube also bounds this difference by $2\epsilon$.

Apply It\^o's formula to $f(\overline Y)$ on each cell, where the
sensor action has already been chosen. The generator and derivative
bounds used in the R11 proof remain valid for this adapted held action.
Stochastic Fubini yields
\begin{align*}
 \int_{t_j}^{t_{j+1}}
       \{f(\overline Y_s)-f(\overline Y_{t_j})\}\,ds
 ={}&\int_{t_j}^{t_{j+1}}(t_{j+1}-v)
       \mathcal A^{\overline a_j}f(\overline Y_v)\,dv\\
 &+\int_{t_j}^{t_{j+1}}(t_{j+1}-v)
             Df(\overline Y_v)\Sigma\,dW_v.
\end{align*}
The $L^2$ norm of the accumulated first term through $t_k$ is at most
$Kht_k/2$. The stochastic integrals are martingale increments over
disjoint cells. It\^o's isometry bounds their accumulated normalized
$L^2$ norm by $Hh\sqrt{t_k/3}$.

Before representation error, every component of
$\Delta W_j-\sqrt h\operatorname{clip}(\Delta W_j/\sqrt h,-z,z)$
has zero mean and variance at most $hv_{\rm clip}$. Components and cells are
independent. Therefore the normalized second moment of the accumulated
projected clipping error is at most
\[
 \frac1d\operatorname{tr}(\Sigma\Sigma^\top)t_kv_{\rm clip}
       =\overline\sigma^2t_kv_{\rm clip}.
\]
The clipping error and the preceding local-drift martingale need not
be independent: the triangle inequality between their already
accumulated norms is sufficient. The implementation error contributes
at most $k\xi_h$.

Let $E_k=\|\overline Y_{t_k}-Z_k\|_{2,d}$. Subtracting the internal
recursion from the integrated sensor state equation thus gives
\[
 E_k\leq\beta h\sum_{j<k}E_j
       +h\sum_{j<k}\|\overline a_j-a_j\|_{2,d}+b_k.
\]
The recursion in~\eqref{eq:r14sensorrecursion} is equivalent to
$e_k=\beta h\sum_{j<k}e_j+h\sum_{j<k}a_j^\Delta+b_k$.
Induction, using the action bound already proved, establishes
$E_k\leq e_k$ and the asserted action inequalities.

The two physical diffusions have the same Brownian term. Subtracting
their integral equations and applying the $\beta$-Lipschitz production
bound gives
\[
 \|\overline Y_t-Y_t\|_{2,d}
 \leq\beta\int_0^t\|\overline Y_s-Y_s\|_{2,d}\,ds
       +\int_0^t a_{\lfloor s/h\rfloor}^\Delta\,ds.
\]
The integral form of Gronwall's inequality proves the bound by $D(t)$.
Projection away from the common-state direction eliminates common
drift, the common schedule, and common diffusion. Bounded production,
the action tube, and It\^o's isometry give, for both physical states,
$\sup_{t\leq T}\|\Pi Y_t\|_{2,d}\leq Q_T$.

It remains to compare payoffs. Subtract the two exact identities
\eqref{eq:r11gain}; the common schedule terms cancel. Production
contributes at most $\beta\int_0^TW(t)D(t)dt$.
For the consumption deficit $R_t$ in~\eqref{eq:r11deficit}, its gradient
in the normalized inner product has components
\[
 -\frac1{m_i}+w(t)+\zeta\bar m
  =\frac1{m_0(t)}-\frac1{m_i}
       +\zeta\{\bar m-m_0(t)\}.
\]
Every point on the segment between the two held actions remains in
their common coordinate tube at $t_k$. Its distance from $m_0(t)$ is
at most $r_k(t)$, and its coordinates are at least $m_{-,k}$. Thus the
normalized gradient norm is bounded by $C_k(t)$. The mean-value
inequality and Cauchy--Schwarz bound the absolute expected deficit
difference by $C_k(t)a_k^\Delta$.

Finally,
\[
 \left|\E\{\|\Pi\overline Y_T\|_d^2-\|\Pi Y_T\|_d^2\}\right|
 \leq\|\overline Y_T-Y_T\|_{2,d}
       \{\|\Pi\overline Y_T\|_{2,d}+\|\Pi Y_T\|_{2,d}\}
 \leq2Q_TD(T).
\]
Adding these three contributions proves
\eqref{eq:r14sensorpayoff}. Adding and subtracting the parent's payoff
on the existing confidence event proves the interval transfer.

\paragraph*{An outward finite computation.}
One need not infer a scalar integral error from a refinement plot.
The function $D$ is nondecreasing and obeys the exact cell recursion
\[
 D(t_{k+1})=e^{\beta h}D(t_k)
                +a_k^\Delta\int_0^h e^{\beta s}\,ds.
\]
An interval enclosure of $m_0$ on each cell gives an upper bound
$C_k^+$ for $C_k(t)$. With the already enclosed weights $A_k,B_k$,
a computable upper allowance is
\[
 \beta\sum_k B_kD(t_{k+1})
       +\sum_k A_kC_k^+a_k^\Delta
       +2\chi e^{-\rho T}Q_TD(T),
\]
evaluated outward. In the forcing recurrence, use nondecreasing saved
upper endpoints $b_k^+\geq b_k$ and the differences of those exact
saved scalars. This preserves the telescoping majorant. The network
and recurrence arithmetic assumptions are the same conditional
binary64 assumptions used for the protected parent.
